% Generated draft: requires mathematical and source review.
\documentclass[11pt]{article}
\usepackage[margin=25mm]{geometry}
\usepackage{fontspec}
\setmainfont{Latin Modern Roman}
\usepackage{amsmath,amssymb,mathtools,unicode-math}
\setmathfont{Latin Modern Math}
\usepackage{enumitem,needspace,xurl,hyperref}
\setlength{\parindent}{0pt}
\setlength{\parskip}{5pt}
\setlength{\emergencystretch}{3em}
\newcommand{\N}{\mathbb{N}}
\newcommand{\Z}{\mathbb{Z}}
\newcommand{\Q}{\mathbb{Q}}
\newcommand{\R}{\mathbb{R}}
\newcommand{\C}{\mathbb{C}}
\newcommand{\F}{\mathbb{F}}
\newcommand{\A}{\mathbb{A}}
\newcommand{\PP}{\mathbb{P}}
\newcommand{\E}{\mathbb{E}}
\newcommand{\eps}{\varepsilon}
\newcommand{\dd}{\,\mathrm{d}}
\begin{document}
\section*{UnsolvedMath ID: 800384}
\textbf{Source problem code:} catalogue entry
\subsection{Definitions and mathematical statement}
Write \(G_K\) for the absolute Galois group of a field \(K\). Whenever a map \(f\) and a basepoint \(P\) determine a rooted inverse-image tree \(T_{f,P}\) with a natural \(G_K\)-action fixing the root, let \(\operatorname{Aut}(T_{f,P})\) denote the group of all rooted-tree automorphisms. The arboreal representation and its Odoni index are
\[
\rho^{\mathrm{arb}}_{f,P}\colon G_K\longrightarrow\operatorname{Aut}(T_{f,P}),
\qquad
\iota_K(f,P):=\bigl[\operatorname{Aut}(T_{f,P}):\operatorname{im}(\rho^{\mathrm{arb}}_{f,P})\bigr].
\]
Here an inverse-image tree has vertices at level \(n\) corresponding to solutions of \(f^n(Q)=P\), with edges induced by \(f\); occurrences at different levels are treated as different vertices. The ambient space, permitted maps and basepoints, and hypotheses ensuring the intended tree and Galois action are not specified in the supplied record. Thus these definitions and the questions below are conditional on the intended source framework. In particular, no regularity or branching degree of the tree is assumed here.

\begin{enumerate}
\item[(a)] Determine precisely which admissible triples \((K,f,P)\) satisfy
\[
\iota_K(f,P)=1,
\]
 equivalently, for which triples \(\rho^{\mathrm{arb}}_{f,P}\) is surjective onto \(\operatorname{Aut}(T_{f,P})\).
\item[(b)] Determine precisely which admissible triples \((K,f,P)\) satisfy
\[
\iota_K(f,P)<\infty,
\]
 equivalently, for which triples the image of \(\rho^{\mathrm{arb}}_{f,P}\) has finite index in the full rooted-tree automorphism group.
\end{enumerate}
Here admissibility refers to the source's intended class of fields, maps and basepoints, which cannot be recovered completely from the supplied excerpt.
\subsection{Short English statement}
Classify the fields, maps and basepoints for which the Galois action on the tree of iterated preimages is surjective, and classify those for which its image has finite index. The supplied excerpt leaves the admissible objects and tree hypotheses unspecified, so the formulation remains conditional on that missing context.
\subsection{Sources}
\begin{itemize}
\item Local UnsolvedMath registry, permanent ID \(800384\), source problem code \texttt{catalogue entry}, titled ``Open Problems in Arithmetic Dynamics --- Question 7.'' The supplied registry labels the problem open; this is a supplied status, not an independently verified literature assessment.
\item Benedetto, DeMarco, Ingram, Jones et al., \emph{Current Trends and Open Problems in Arithmetic Dynamics} (2018), Question 7. Supplied source URL: \url{https://arxiv.org/abs/1806.04980}. The registry attributes the proposal to these authors in 2018.
\item The supplied provenance records access on 2026-10-16 and extraction by \texttt{arxiv\_latex}. These are supplied provenance details; no source visit or independent verification is claimed.
\end{itemize}
\end{document}
